%% feature/sec-introduction
\section{Introduction}
Reasoning models are good because they think out loud, and expensive for exactly
the same reason. A model that solves a hard problem by writing thousands of
intermediate tokens pays for every one of them in latency and cost, and re-derives
the same sub-procedures---the same case splits, the same standard substitutions---on
problem after problem. The obvious fix is to hand the model the result of that
deliberation up front: a short, distilled insight that says, in effect, ``you have
seen this shape before, here is the trick.'' Inserted once, such an insight already
works. It preserves accuracy while cutting generated tokens, and it does so without
fine-tuning or long demonstrations.

But once is a strange place to stop. A distilled insight is short, self-contained,
and already validated as helpful---which makes it precisely the kind of thing one
would expect to gain from emphasis. And there is direct precedent for emphasis
through sheer repetition: simply duplicating a prompt, with no new information added,
measurably improves non-reasoning models. So a natural lever sits in plain sight. If
showing the insight once helps, does showing it twice, three times, five times help
more? Does it keep the model on the rails for longer, or compress its reasoning
further?

Unfortunately, the literature points in two directions at once, and the conflict is
the whole reason this is worth testing rather than assuming. Repetition helps
non-reasoning models---but the same intervention is nearly inert once reasoning is
switched on, because reasoning traces already restate the prompt on their own. A
mechanistic account of why content-light repetition helps predicts a benefit for
uniform, low-attention insertions, yet predicts the opposite---active
harm---for diverse, structured, semantically heavy insertions. A distilled insight is
exactly the latter kind of object. And feeding a reasoning model raw reasoning
content from a similar problem does not merely fail to help; it actively degrades
accuracy, the more of it you add. Repeating a fixed, meaningful insight could
therefore behave like a benign reinforcing signal or like an over-constraining
perturbation that inflates the prompt faster than it compresses the output. The sign
is genuinely unknown.
% Preamble: \usepackage{tikz}  \usetikzlibrary{arrows.meta}
\begin{figure*}[t]
\centering
\resizebox{\textwidth}{!}{%
\begin{tikzpicture}[
    font=\small,
    >={Latex[length=2mm]},
    box/.style={draw=#1!70!black, fill=#1!12, rounded corners=2pt,
                align=center, inner sep=5pt, minimum height=11mm},
    card/.style={draw=violet!70!black, fill=violet!10, rounded corners=2pt,
                 align=center, inner sep=2pt, minimum width=30mm,
                 minimum height=5mm, font=\footnotesize},
    out/.style={draw=#1!70!black, fill=#1!10, rounded corners=2pt,
                align=center, inner sep=4pt, minimum width=40mm,
                minimum height=13mm, font=\footnotesize}
]

% ===== Left: the established single-insertion fact =====
\node[font=\small\bfseries] at (0, 2.55) {Established ($k{=}1$)};
\draw[draw=black!40, dashed, rounded corners=3pt] (-2.4, 2.15) rectangle (2.4, -1.55);
\node[card, draw=black!45, fill=black!6] (q1)  at (0, 1.55) {Problem $q_i$};
\node[card] (i1) at (0, 0.80) {Insight $e_i$};
\node[font=\footnotesize, text=teal!60!black] at (0, -0.05) {$\Downarrow$ one copy helps};
\node[out=teal, minimum width=34mm] (base) at (0, -0.95)
    {accuracy preserved,\\ fewer generated tokens};

% ===== Center: the intervention and the open question =====
\node[font=\small\bfseries] at (6.7, 2.55) {This paper ($k>1$)};
\draw[draw=violet!55, dashed, rounded corners=3pt] (4.3, 2.15) rectangle (9.1, -1.55);
\node[card, draw=black!45, fill=black!6] (q2) at (6.7, 1.55) {Problem $q_i$};
\node[card] (ia) at (6.7, 0.85) {Insight $e_i$};
\node[card] (ib) at (6.7, 0.25) {Insight $e_i$};
\node[font=\footnotesize] at (6.7, -0.30) {$\vdots$ \;\; (same card, $k\times$)};
\node[font=\footnotesize, text=violet!60!black] at (6.7, -1.05)
    {repeat the \emph{identical} insight};

\node[font=\Large\bfseries, text=black!70] at (11.0, 0.30) {?};

% ===== Right: the three competing outcomes =====
\node[out=teal,   minimum width=46mm] (o1) at (15.6, 1.55)
    {\textbf{Help.} repetition reinforces\\ the hint: higher accuracy or\\ fewer tokens};
\node[out=gray, fill=black!6, minimum width=46mm] (o2) at (15.6, -0.05)
    {\textbf{Saturate.} extra copies add\\ nothing beyond the first};
\node[out=orange,  minimum width=46mm] (o3) at (15.6, -1.65)
    {\textbf{Hurt.} over-constrains and\\ inflates the prompt faster\\ than it compresses output};

% ===== Arrows =====
\draw[->] (2.45, -0.95) .. controls (3.4, -0.95) and (3.4, 0.30) .. (4.25, 0.30);
\draw[->] (9.15, 0.30) -- (10.6, 0.30);
\draw[->] (11.45, 0.30) -- (12.4, 1.55) -- (o1.west);
\draw[->] (11.45, 0.30) -- (12.4, -0.05) -- (o2.west);
\draw[->] (11.45, 0.30) -- (12.4, -1.65) -- (o3.west);

% ===== Footnote =====
\node[font=\footnotesize, text=black!60] at (7.8, -2.6)
    {Only the repetition count $k$ changes; the insight, question, model, and decoding are held fixed.};
\end{tikzpicture}}
\caption{The question this paper isolates. A distilled reasoning insight inserted
\emph{once} ($k{=}1$, left) already preserves accuracy while cutting generated
tokens. We repeat the \emph{same} insight $k$ times before the problem (center),
changing nothing else, and ask which of three outcomes follows (right): repetition
helps, saturates, or hurts. \textbf{Takeaway:} the literature predicts all three for
this cell, so the sign is genuinely open---which is exactly why it is worth a
controlled test.}
\label{fig:teaser}
\end{figure*}

We isolate this one effect cleanly. Each problem comes with a pre-built, independently
validated insight---a reasoning skill card already shown to help at a single
insertion---so the question ``does the insight help'' is settled in advance and held
fixed. We then repeat that identical card $k$ times before the problem and change
\emph{nothing else}: same question, same gold answer, same insight text, same prompt
template, same model, same decoding. Repetition count is the only moving part, so any
difference we measure cannot be blamed on a better retriever, a sharper extractor, or
a luckier insight. Against this we run matched baselines that repeat the bare problem,
repeat a raw reasoning trace, and pad with filler, all at the same copy counts, and we
score every condition on two axes that matter for deployment: accuracy, and the total
tokens it costs to get there.

This is a pre-results manuscript: we specify the experiment and the evidence it must
produce, not the answer. Our contributions are the design itself.

\begin{itemize}
  \item \textbf{A new question.} We formalize \emph{insight repetition}---repeating a
        fixed, pre-validated reasoning insight, rather than a raw prompt---as the
        untested intersection of prompt repetition and distilled-insight transfer.
  \item \textbf{A confound-free design.} A copy-count-matched protocol that varies only
        the repetition count $k$, holding the insight and everything else fixed, with
        matched baselines (bare-problem, raw-trace, and padding repetition) at every
        $k$.
  \item \textbf{Efficiency as a first-class axis.} We report total generated tokens,
        not only visible or hidden reasoning tokens, so the input inflation repetition
        causes is made explicit and traded off against any accuracy change.
  \item \textbf{Generality across models.} The protocol spans more than one reasoning
        model family, separating model-specific quirks from effects that generalize.
\end{itemize}
